\documentclass[11pt, a4paper]{article}
\usepackage[utf8]{inputenc}
\usepackage{amsmath, amssymb}
\usepackage{graphicx}
\usepackage[margin=1in]{geometry}
\usepackage{tcolorbox}
\usepackage{fancyhdr}

\definecolor{UNIBBlue}{RGB}{0, 51, 102}

\pagestyle{fancy}
\fancyhf{}
\rhead{\textbf{Optimisasi untuk Teknik Elektro}}
\lhead{Problem Set - Minggu 4}
\cfoot{\thepage}

\begin{document}

\begin{center}
    \Large\color{UNIBBlue}\textbf{PROBLEM SET}\\
    \large\textbf{Minggu 4: Analisis Sensitivitas}
\end{center}

\vspace{0.5cm}

\textbf{Instruksi:} Kerjakan soal-soal di bawah ini dengan menyertakan langkah-langkah analitis dan dasar teori analisis sensitivitas.

\section*{Soal 1: Terminologi (C1 - C2)}
1. (C1) Apakah perbedaan antara Range of Optimality dan Range of Feasibility? \\
2. (C2) Mengapa analisis sensitivitas sangat penting bagi seorang insinyur perencana (\textit{planning engineer}) sistem telekomunikasi ketika memprediksi kebutuhan infrastruktur jaringan ke depan?

\section*{Soal 2: Evaluasi Skenario (C4)}
Diberikan laporan sensitivitas untuk maksimasi kinerja produksi panel surya Tipe A ($x_1$) dan Tipe B ($x_2$):

\begin{table}[h]
\centering
\begin{tabular}{|l|c|c|c|c|}
\hline
\textbf{Produk} & \textbf{Final Value} & \textbf{Obj. Coef. (\$)} & \textbf{Allowable Increase (\$)} & \textbf{Allowable Decrease (\$)} \\ \hline
Tipe A ($x_1$) & 120 & 40 & 10 & 25 \\ \hline
Tipe B ($x_2$) & 80  & 30 & 50 & 6  \\ \hline
\end{tabular}
\end{table}

(a) Tentukan rentang optimalitas untuk koefisien tujuan dari Tipe A ($c_1$) dan Tipe B ($c_2$). \\
(b) Jika karena kelangkaan material tipe B, harga jual (objektif) tipe B turun menjadi \$25, apakah konfigurasi optimal produksi berubah?

\section*{Soal 3: The 100\% Rule (C5)}
Dengan menggunakan data dari Soal 2, asumsikan terjadi perubahan secara bersamaan: profit Tipe A meningkat sebesar \$5 (karena subsidi hijau pemerintah), dan profit Tipe B menurun sebesar \$4 (karena kalah saing). 

Gunakan \textit{The 100\% Rule} untuk menentukan apakah solusi optimal saat ini terjamin tetap optimal, atau ada kemungkinan berubah. Tunjukkan perhitungan persentasenya secara lengkap.

\section*{Soal 4: Sintesis Pengambilan Keputusan (C6)}
Sebuah analisis alokasi spektrum frekuensi menghasilkan \textit{shadow price} sebesar 5.0 (Mb/s) per MHz untuk alokasi band 5GHz, dengan \textit{allowable increase} sebesar 20 MHz.
Saat ini regulator menawarkan izin blok tambahan 15 MHz di band 5GHz dengan biaya sewa tahunan yang setara jika dikonversi menjadi denda performansi senilai kehilangan throughput total 60 Mb/s.

(C6) Sebagai Chief Technology Officer (CTO), apakah Anda akan membeli blok 15 MHz tersebut? Sertakan argumen matematis berdasarkan perbandingan \textit{shadow price} dan biaya margin, serta nyatakan validitasnya melalui \textit{Range of Feasibility}.

\end{document}
